\documentclass[UTF8,11pt,a4paper,fontset=fandol]{ctexart}
\usepackage[margin=23mm,headheight=15pt]{geometry}
\usepackage{amsmath,amssymb,booktabs,tabularx,longtable,array,graphicx,xcolor}
\usepackage{tikz}
\usetikzlibrary{arrows.meta,positioning}
\usepackage{hyperref,fancyhdr}
\hypersetup{colorlinks=true,linkcolor=blue!45!black,citecolor=blue!45!black,urlcolor=blue!45!black,pdftitle={机器人连续学习中的能力保留},pdfauthor={LI HaoRan}}
\setlength{\parindent}{2em}
\setlength{\parskip}{3pt}
\setlength{\emergencystretch}{2em}
\renewcommand{\arraystretch}{1.2}
\pagestyle{fancy}
\fancyhf{}
\fancyhead[L]{\small 研究报告 1}
\fancyhead[R]{\small 视觉—语言—动作模型的连续微调}
\fancyfoot[C]{\thepage}
\newcolumntype{Y}{>{\raggedright\arraybackslash}X}
\title{\vspace{-1.5em}\textbf{机器人连续学习中的能力保留}\\[0.4em]\large 视觉—语言—动作模型微调的遗忘机制、抑制方法与评价边界}
\author{LI HaoRan}
\date{2026年11月27日}
\begin{document}
\maketitle
\begin{abstract}
视觉—语言—动作（VLA）模型能够借助预训练知识学习新的机器人操作，但适应新任务的参数更新也可能破坏已有技能。本文围绕这一问题，评述十篇于2026年1月之后首次公开的研究，比较经验回放、保守更新、表征约束、指令锚定与强化微调的作用机制。现有证据支持有条件的能力保留：预训练能够提高连续适应的效率，但效果仍受学习目标、历史信息覆盖、动作归一化和测试分布影响。分布内任务成功、旧指令响应和分布外鲁棒性是不同的评价对象。本文据此区分各研究的实验证据与解释范围，并提出同时评价旧任务保留、新任务习得、分布外鲁棒性及资源成本的实验设计。本文采用文献证据综合方法，比较数据均引自原始研究。
\end{abstract}

\section{研究问题与文献范围}
机器人投入使用后，需要不断学习新的操作，又要保留已经掌握的技能。例如，学会按照新指令把碗放到柜子顶部以后，仍应能够按照旧指令把碗放到炉台上。即使机器人保留了视觉识别能力，并且机械结构能够执行两种运动，训练后也可能对熟悉场景产生错误的目标选择。因此，连续学习中的失败并不必然等同于所有知识被抹除。

本文关注的技术问题是：\emph{顺序微调过程中，数据覆盖、参数更新和多模态对齐如何导致旧技能退化，哪些干预能够缓解这种退化？}主要评价对象是已经学习过的机器人操作行为。基础视觉语言模型的通用问答能力属于相关但不同的能力，保留通用语义能力不能直接证明旧操作技能得到保留。

文献纳入范围为2026年1月之后首次公开、截至2026年11月27日可获得的十篇研究。检索结合主题词、相关工作及官方arXiv或会议记录，优先纳入直接研究连续VLA学习或受控旧能力保持的论文。文献的同行评审状态并不相同，arXiv编号本身不能证明已经通过同行评审。其中Liu等的研究已有ICML 2026会议论文集记录；近期预印本因与问题直接相关而纳入，并保留相应证据限制。本文为聚焦问题的批判性综述，不主张穷尽全部相关研究。

本文重点考虑30B以内的模型，核心比较多使用约3B--7B参数量级的VLA，包括OpenVLA、$\pi_0$、$\pi_{0.5}$和GR00T系列。LifeLong-RFT中NORA-Long的精确总参数量未在所审阅材料中明确，因此不将全部系统均低于30B作为已确认结论。可部署性还取决于训练状态、图像token数、动作预测长度、在线采样以及教师模型成本，不能仅凭参数量判断。

\section{从策略学习理解遗忘}
\subsection{VLA模型学习的对象}
VLA策略接收图像观察$o$、语言指令$\ell$及可能的机器人状态，输出动作或一段连续动作$a$：
\begin{equation}
a\sim\pi_\theta(a\mid o,\ell).
\end{equation}
视觉编码器提取场景特征，多模态主干结合语言，动作模块生成控制命令。自回归策略将动作表示为token并依次预测；回归头直接输出连续动作；流匹配头则学习速度场，通过积分把噪声变换为动作。三者的输出形式和损失函数不同，保持历史响应时可施加的约束也不同。

架构决定保持约束作用的对象与成本：自回归策略输出动作token分布，回归头输出连续动作块，流匹配头则预测含噪动作轨迹上的速度场。XIA在相同条件下匹配架构特定的师生响应，包括流速度；ConSFT缩放流匹配监督更新，不增加参考网络；Info-VLA冻结视觉语言编码器、更新动作侧组件，因此编码器权重不变仍不能证明下游控制得到保留。这些区别解释约束应施加在哪里，尚不足以证明某类输出架构天然更少遗忘\cite{xia,consft,info}。

在连续学习中，任务数据$D_1,\ldots,D_K$依次到达，学习第$k$个任务时，完整历史数据往往不再可用。普通监督微调直接拟合当前专家示范；经验回放（experience replay，ER）额外保存少量历史数据$B_{k-1}$，其基本思想可表示为
\begin{equation}
\mathcal L_k=\mathbb E_{D_k}[\ell_{\rm BC}]+\lambda\mathbb E_{B_{k-1}}[\ell_{\rm BC}].
\label{eq:er-zh}
\end{equation}
这里$\ell_{\rm BC}$为行为克隆损失。具体实现可以在一个批次中混合新旧样本，也可以交替使用完整批次。缓存容量与回放频率是两个控制量，保留更多帧不等于训练时更频繁地使用这些帧。

\subsection{新任务更新为何会损害旧任务}
对旧任务损失作局部展开，可以得到
\begin{equation}
\Delta L_{\rm old}\approx\nabla L_{\rm old}^{\top}\Delta\theta+
\tfrac12\Delta\theta^{\top}H_{\rm old}\Delta\theta.
\label{eq:taylor-zh}
\end{equation}
远离旧任务最优点时，一阶项仍然重要；接近最优点时，部分研究使用Fisher矩阵近似曲率，分析新任务梯度$g$沿敏感方向的量$g^{\top}Fg$\cite{hu_rl,consft}。这类分析描述的是局部损失敏感性，并不能精确预测闭环操作成功率。很小的参数变化也可能使抓取动作跨过接触或稳定性阈值。

低秩适应（LoRA）把权重写作$W=W_0+BA$，其中秩$r$远小于原矩阵维度。它通过减少可训练参数、约束更新方向来降低微调成本，但一个反复复用的适配器也会遗忘。LoRA没有自动为每个任务分配互不干扰的参数，其效果仍依赖训练目标、秩、主干和训练过程。

\subsection{同时衡量保留与习得}
令$R_{k,i}$表示完成前$k$个任务训练后，在任务$i$上的成功率。为明确评价对象，定义以下最终检查点汇总指标：
\begin{align}
A_K&=\frac1K\sum_{i=1}^{K}R_{K,i},\\
F_K&=\frac1{K-1}\sum_{i=1}^{K-1}(R_{i,i}-R_{K,i}).
\end{align}
$F_K$是从任务习得时到最终检查点的带符号遗忘量，负值表示后续训练带来改善。此处用于统一说明评价含义，各原论文未必采用同一公式。有些方法对所有后续检查点取平均，有些只比较最终检查点；后向迁移与负后向迁移的符号约定也可能相反。未重新计算以前，不应将这些数值直接合并排名。

遗忘量还需与习得表现联合解释：模型若未充分学会某任务，后续的性能下降也可能很小。因此，需要同时报告任务刚学会时的表现、最终表现和完整的任务—检查点矩阵。

\section{十篇研究构成的机制证据图谱}
表\ref{tab:map-zh}按各研究检验的机制组织文献。各研究的任务、资源和评价口径不同，不宜据此进行统一性能排名。
\begin{table}[htbp]
\centering\small
\begin{tabularx}{\textwidth}{@{}p{27mm}YY@{}}
\toprule
研究 & 主要机制或干预 & 解释时必须保留的条件\\
\midrule
Liu等\cite{liu_er} & 预训练、回放、组件替换与技能恢复 & 以仿真为主，初始能力与回放量影响结果\\
Hu等\cite{hu_rl} & 预训练、LoRA与在策略强化学习 & 需要交互式强化信号，不等同离线模仿\\
Zhu等\cite{zhu_real} & 真实任务流、回放调度与动作归一化 & 分步得分与固定机器人平台\\
ConSFT\cite{consft} & 按样本损失保守缩放梯度 & 主要考察目标适应与原能力保持\\
PHASER\cite{phaser} & 分阶段记忆分配与干扰感知路由 & 阶段标签来源及近似匹配的缓存预算\\
XIA\cite{xia} & 历史指令条件下的响应锚定 & 增加教师查询与临时锚点存储\\
OOD分析\cite{ood} & 旧任务分布外表现与表征探针 & 单模型、单训练种子及相关性证据\\
RECALL\cite{recall} & 不确定性感知的恢复数据采集 & 采集质量与回放整合存在耦合\\
Info-VLA\cite{info} & 跨模态信息约束与教师引导 & 对齐估计是行为保持的代理目标\\
LifeLong-RFT\cite{lifelong} & 基于专家参照奖励的强化微调 & 奖励设计和历史示范覆盖\\
\bottomrule
\end{tabularx}
\caption{十篇文献的机制与证据边界。}
\label{tab:map-zh}
\end{table}

\section{预训练、学习目标与真实执行的共同作用}
\subsection{预训练提高回放效率，但不能据此推断无需防遗忘}
Liu等将预训练VLA与较小策略比较，并在相同架构内部控制初始化水平\cite{liu_er}。这些实验支持一个有条件的结论：较强的预训练先验能够使少量历史数据发挥更大作用。相同架构内部的对照，比跨架构直接将差异归因于参数量更有解释力。

该研究进一步通过组件替换及恢复训练，检验技能退化是否意味着知识完全消失。其结果支持有用表征的部分保留。表4中$\pi_0$在四个任务集上的恢复步数与最初学习步数之比分别为0.066、0.105、0.067和0.062，即约6.2\%--10.5\%。这些结果说明恢复成本较低，但并非所有任务集均低于初次训练步数的10\%。恢复仍依赖旧任务数据及额外训练，并不等同于部署策略自主恢复。

\subsection{训练目标会改变持续适应的难度}
Hu等研究预训练VLA、LoRA和在策略强化学习共同作用下的连续适应\cite{hu_rl}。在LIBERO-Spatial上，其表1报告最终平均成功率为$81.2\pm0.4$\%，负后向迁移（NBT）为$0.3\pm0.5$个百分点。将强化学习目标替换为SFT后，消融结果为$29.9\pm2.3$\%平均成功率和$78.7\pm1.9$个百分点NBT。该对照说明，在所测条件下，初始化、适应方式与优化目标存在重要交互。

这里的在策略强化学习由当前策略采样动作，并通过交互获得评价；离线监督模仿则拟合外部专家示范。两者的训练状态分布、监督信号和资源成本不同，不能将强化学习结果直接推广到离线SFT。高维空间中可能存在较少干扰的方向，也不意味着训练梯度实际随机或必然安全。Fisher及更新方向分析是对特定模型的证据，不是全部大模型不会遗忘的普遍定理。

\subsection{真实机器人暴露动作表示层面的失败}
Zhu等建立了包含单臂和双臂操作的十项真实任务流\cite{zhu_real}。在五任务同构流中，其表I给出的单臂最终分步平均得分为顺序微调31.4、ER 97.2；双臂分别为38.3和97.2。这些数值是在0--100尺度上的任务进度得分，并非二元任务成功率，不能和LIBERO成功率直接比较。

动作归一化实验提供了更具体的机制线索。改变归一化统计量会改变动作目标的数值坐标系；如果同一物理动作在不同阶段对应不一致的数值目标，历史回放便难以保持原映射。在其测试任务中，固定首任务统计量在覆盖足够动作范围时有效；如果未来动作超出该范围，就不能保证这一策略仍适用。异构任务流还显示双臂技能需要非对称的回放调度。因此，是否采用ER只是持续部署设计中的一部分。

\section{抑制遗忘的方法及其代价}
\subsection{使缓存覆盖关键阶段和易混指令}
PHASER将任务内部的阶段分配与跨任务回放路由分别处理\cite{phaser}。均匀抽帧容易遗漏持续时间短、却决定整项操作能否成功的抓取阶段。按阶段提供支持旨在缓解这一瓶颈；多模态路由则优先处理视觉和语言相近、所需动作却不同的上下文。其表1中，OpenVLA-OFT在LIBERO-Long上的最终平均成功率由ER的54.6\%提高至PHASER的85.8\%，差为31.2个百分点。

这一比较允许约$\pm10$\%的缓存量变化，因此“预算匹配”应理解为近似匹配。主要实验使用人工阶段标注，Auto-PC作为自动阶段识别方案另行评价。结果支持所测条件下的阶段感知记忆分配，比较方法成本时仍需计入阶段识别和预处理开销。

XIA针对另一类覆盖空缺\cite{xia}。即使增加缓存中样本的权重，也不能直接监督未在缓存中共同出现的“当前观察—旧指令”组合。XIA将当前观察与历史指令组合，匹配前一策略的响应，可抽象为
\begin{equation}
\mathcal L_{\rm XIA}=\mathcal L_{\rm ER}+\alpha\mathbb E\!\left[
d\left(r_\theta(q,\ell^-),\operatorname{sg}[r_{\theta^-}(q,\ell^-)]\right)\right].
\end{equation}
其中$q$为响应查询：直接动作预测使用观察，流策略则使用观察、含噪动作块及流时间；$\ell^-$为历史指令，$r$为策略响应，$d$为响应间的差异度量，$\alpha$为约束权重，$\theta^-$为冻结的上一阶段策略，$\operatorname{sg}$表示停止教师梯度。历史条件下的目标来自教师响应，而非将新任务专家动作重新标注为旧任务标签。XIA报告OpenVLA-OFT在三个LIBERO任务集上的平均最终成功率由ER的72.5\%提高到92.6\%。它保持了专家帧预算，同时增加教师计算和临时锚点数据。教师可能犯错，旧指令在新观察中也可能无法完成，因而响应保持的收益仍需结合这些条件理解。

\subsection{以保守更新保护历史能力}
ConSFT按样本损失对监督梯度作指数缩放\cite{consft}：
\begin{equation}
\mathcal J=\operatorname{sg}\!\left[e^{-\mathcal L/\tau}\right]\mathcal L,
\qquad g_{\rm ConSFT}=e^{-\mathcal L/\tau}g_{\rm SFT}.
\end{equation}
高损失样本的更新贡献更小；在局部曲率矩阵固定的近似下，相应二次敏感性按权重的平方缩放。该推导并不证明长时闭环行为具有统一上界，也不保证陌生但必要的技能能够顺利习得。论文指出，不熟悉控制原语的学习可能变慢，且未给出多步流轨迹的形式化保持界限。

上式描述基本权重机制。原论文附录C.1中的实际权重还包含损失尺度校准$\kappa$、最小权重$\omega_{\min}$和随训练步数变化的温度：
\begin{equation}
\omega_t=\max\!\left(\omega_{\min},\exp\left[-\kappa\mathcal L/\tau_t\right]\right).
\end{equation}
最小权重避免高方差有效样本的梯度被完全抑制，温度升高则逐渐减弱约束。因此，实际训练使用带下限和温度调度的权重，不能仅以固定温度的基本公式描述。

其表2中，$\pi_0$旧任务集平均成功率由SFT的0.09提高至ConSFT的0.34，同时目标任务成功率均为0.90。这一比较中的保留优势为25个百分点，但相较基座的0.49仍存在明显能力损失。该实验主要研究窄目标适应时的基础能力保持，不能等同于回放研究中的十任务长期顺序学习。

\subsection{表征约束、主动采集与强化微调}
Info-VLA将回放、教师引导的表征对齐及信息论一致性目标结合\cite{info}。作者冻结视觉语言编码器，更新包含融合模块在内的动作侧组件。表1中，LIBERO-Long最终成功率由ER的66.6\%提高到组合方法的73.6\%。但表3同时报告ER的NBT为$-4.3$，组合方法为$-1.0$，并标注数值越小越好。因此，最终成功率提升并不代表每项保留指标都改善。互信息估计及表征梯度路径还应结合代码核查，代理目标本身不足以确立信息保持保证。

RECALL使用基于不确定性的求助预测器发现需要恢复示范的状态\cite{recall}。在线方案从首次触发求助的状态开始采集；离线替代方案从多个被标记的仿真状态采集。在配合广泛回放的比较中，由Strong INSIGHT引导的在线方案总体成功率为72.4\%，匹配的被动采集方案为60.2\%。Strong INSIGHT指求助预测器的训练标签设置，不表示它在所有条件下都更优。在线与离线采集分别为72.4\%和68.4\%，但论文给出的$p=0.1659$未提供在线方案统计上更优的充分证据。主实验中仅恢复数据与回放方案还使用了不同的动作归一化约定，因此全部差距不能单独归因于回放。该研究围绕LIBERO-10上的恢复更新周期，证据范围窄于无限期积累任务的连续学习主张。

LifeLong-RFT根据专家参照对采样动作块评分，奖励包含token一致性、解码后位姿与夹爪对齐、输出格式合法性\cite{lifelong}。裁剪的组相对更新及参考策略正则化提供了另一种适应路径。奖励计算不需要在线执行环境，但连续学习协议仍保留每个历史任务五条示范。因此它既不同于无回放适应，也不同于通过环境交互获取奖励的强化学习。

其表IV中四个任务集学习曲线下面积（AUC）差值的平均值约为21.9个百分点，这既不是相对增长21.9\%，也不是最终检查点成功率的差值。专家参照奖励还可能惩罚与示范轨迹不同、却物理上同样可行的动作。以上三类方法表明，数据访问条件、可训练模块及优化约束需要共同审视。

\section{如何理解看似冲突的研究结论}
\subsection{分布内保留不等于同一旧任务在扰动下仍然可靠}
OOD研究将ER训练的$\pi_0$放到LIBERO-Plus扰动下重新评价\cite{ood}。模型分布内保留较好时，旧任务分布外表现仍可能随连续训练退化。该研究区分初始域偏移造成的损失与后续任务训练带来的进一步下降。单纯观察到OOD得分低不能证明连续遗忘，需要在相同扰动条件下沿检查点追踪同一旧任务。

线性探针发现，在偏移域内部重新拟合读出器，可以提取原分布读出器未能有效利用的任务信息。这与表征失配的解释一致，但尚不是唯一因果机制的证明。单模型、单训练种子及相关性分析限制了外推范围。由此得到的可靠结论是：在这一设置下，仅凭分布内保留不能证明部署鲁棒性，需要进一步受控检验。

\subsection{比较方法时还要比较可用信息和成本}
看似矛盾的结果，往往来自三个层面的差别。首先，离线模仿与在线强化学习的监督来源和状态分布不同。其次，回放可以按帧或完整轨迹计量，可以固定每任务容量或固定总容量，也可以改变新旧数据混合频率。固定每任务容量仍会使总存储随任务数增长。最后，各研究在预训练重叠、任务多样性、动作统计、测试分布和评分定义上并不相同。

因此，同样叫ER并不意味着同样的实验；新增方法优于某项ER基线时，还要核对具体缓存和训练配置。专家帧预算相同也不代表总计算相同，教师查询、临时存储和交互次数都可能额外增加。训练步数一致同样不能证明FLOPs、耗时或机器人交互成本一致。

\section{能够区分机制的后续实验}
值得进一步检验的问题是：\emph{在动作统计受控、历史信息预算明确时，阶段感知和指令感知回放能否同时保留旧技能及其分布外鲁棒性？}以下提出能够检验该问题的对照设计。

选择一个公开可用、参数量不超过30B的VLA并固定初始化，披露预训练任务重叠。比较顺序SFT、ER、阶段感知回放、ER加指令锚定和保守更新基线。交互式强化学习作为独立实验组，单列采样与环境成本。在适用范围内统一当前任务示范、优化预算、任务顺序及评价种子，并分别报告固定每任务缓存和固定总缓存两种制度。

每完成一个任务，都评价全部已学习任务及预设的相机、光照、物体布局和指令扰动，保存$R_{k,i,p}$，其中$k$为检查点、$i$为任务、$p$为扰动条件。增加目标混淆矩阵，区分完全未完成动作与执行了错误指令。联合报告习得、最终成功、遗忘及独立训练种子之间的不确定性，同时记录存储字节数、教师查询、显存与训练或采样时间。

为区分机制，在固定缓存总量下消融阶段分配；在固定教师查询次数下，将历史指令锚点替换为当前指令锚点；比较固定动作统计和因果可实施的增量统计；检验读出修复是否真正改善闭环任务成功。仅依赖表征探针或CKA相似度仍属于描述性证据，能够在控制其他条件时修复行为的干预，才提供更强的因果支持。

\begin{figure}[htbp]
\centering
\begin{tikzpicture}[node distance=6mm,every node/.style={font=\small},box/.style={draw=blue!50!black,fill=blue!4,rounded corners,align=center,text width=35mm,minimum height=13mm},>=Latex]
\node[box] (data) {新任务数据\\有界历史信息};
\node[box,right=of data] (update) {覆盖关键条件\\约束更新与对齐};
\node[box,right=of update] (test) {旧技能与新技能\\原分布与扰动分布};
\draw[->,thick] (data)--(update);
\draw[->,thick] (update)--(test);
\end{tikzpicture}
\caption{本文的分析框架。评价保留主张，需要同时核对可用信息、学习干预和行为测试。图中呈现综述的概念结构，不表示实验测量结果。}
\end{figure}

\section{结论}
当前证据支持的是有条件的能力保留。预训练与受约束的适应能够使简单连续学习方法取得良好表现，但历史信息覆盖和动作表示约定决定了哪些行为真正受到保护。阶段感知回放针对时间上的关键瓶颈，指令锚定针对未受到监督的条件响应，保守目标针对破坏性较强的更新；任何单一干预都不足以确立普遍部署可靠性。

关键待解问题是：当环境发生变化、任务流持续延长，而存储与交互成本受到约束时，这些方法能否继续保持旧技能。回答这一问题，需要联合报告保留、习得、鲁棒性与资源。本综述据此将乐观的仿真结果、真实执行失败及OOD退化放回各自实验条件中解释，并保留每项研究实际证明了什么的边界。

\appendix
\section{文献日期、分析版本与发表状态}
纳入日期按首次公开提交计算，不用后续修订日期替代。下表十篇文献均首次公开于2026年1月31日之后，日期及版本依据截至2026年11月27日的官方arXiv记录。会议关联与本文实际使用的预印本版本分别记录；已获接收不意味着所读预印本与最终会议版完全相同。

\begingroup
\footnotesize
\setlength{\tabcolsep}{3.5pt}
\begin{longtable}{@{}p{25mm}p{24mm}p{20mm}p{23mm}p{47mm}@{}}
\caption{文献纳入日期及来源状态。}\label{tab:corpus-zh}\\
\toprule
研究 & arXiv编号 & 首次公开 & 分析版本 & 发表状态依据\\
\midrule
\endfirsthead
\toprule
研究 & arXiv编号 & 首次公开 & 分析版本 & 发表状态依据\\
\midrule
\endhead
\bottomrule
\endfoot
Liu等\cite{liu_er} & 2603.03818 & 2026-03-04 & v2，03-17 & ICML 2026论文集收录，本文分析arXiv v2。\\
Hu等\cite{hu_rl} & 2603.11653 & 2026-03-12 & v3，07-11 & 官方arXiv备注报告RLC 2026接收。\\
Zhu等\cite{zhu_real} & 2605.26820 & 2026-05-26 & v3，08-08 & arXiv预印本，纳入依据为arXiv预印本记录。\\
ConSFT\cite{consft} & 2605.08879 & 2026-05-09 & v2，05-19 & arXiv预印本，纳入依据为arXiv预印本记录。\\
PHASER\cite{phaser} & 2606.03598 & 2026-06-02 & v2，06-03 & arXiv预印本，纳入依据为arXiv预印本记录。\\
XIA\cite{xia} & 2611.31903 & 2026-11-24 & v1，11-24 & arXiv预印本，距综述截止日仅三天。\\
OOD分析\cite{ood} & 2611.31511 & 2026-11-23 & v1，11-23 & arXiv备注报告CoRL 2026 Learn@Deploy工作坊展示，不作为主会论文集声明。\\
RECALL\cite{recall} & 2606.23617 & 2026-06-22 & v1，06-22 & 纳入依据为arXiv预印本记录。\\
Info-VLA\cite{info} & 2603.13335 & 2026-03-06 & v1，03-06 & 作者机构发表清单列为CoRL 2026，分析使用3月预印本。\\
LifeLong-RFT\cite{lifelong} & 2602.10503 & 2026-02-11 & v2，05-16 & 作者项目引用arXiv预印本。\\
\end{longtable}
\endgroup

Liu等的会议记录见\url{https://proceedings.mlr.press/v306/liu26bu.html}，Info-VLA的会议关联见作者机构清单\url{https://milab.westlake.edu.cn/publications.html}。其余状态依据参考文献所链接的官方arXiv记录。仅列出arXiv预印本记录的条目，其期刊或会议接收状态未纳入本文的已确认信息。

质量评价同时考虑研究问题的相关性、基线、消融、不确定性报告及复现细节。六篇仅列为预印本的研究，其同行评审状态未获确认；两篇11月新稿可供独立检验的时间尤其有限。因此，涉及这些结果的综合判断仍需后续独立复现和更广泛评测支持。

核心常见VLA系列的近似规模低于30B，但NORA-Long的精确总参数量尚未确认。主干参数量、可训练参数量与部署策略总量并非同一指标，教师、阶段分割及求助预测等辅助组件也应计入实际成本。

\bibliographystyle{unsrt}
\bibliography{references}
\end{document}
